%!TEX root = ../thesis.tex
\chapter{Conclusion and Outlook}
\label{ch:conclusion}

\section{Conclusion}
The research questions formulated in this study were successfully addressed.
The results indicate that additional image channels, such as infrared (IR), can enhance detection performance, particularly for small objects, whereas spectral indices like NDVI do not necessarily lead to improvements.
The experiments demonstrate that RGB data alone provides a robust foundation for detecting small vehicles, although the inclusion of IR data further increases accuracy.
Moreover, it became evident that the availability of extensive, high-quality training data has a greater influence on model performance than either expanding the input channels with additional spectral information or using a more recent model version.
Overall, the findings confirm that \acrshort{YOLO}v9 is well suited for small object detection and that architectural advancements over earlier versions, such as \acrshort{YOLO}v3, yield measurable gains in both classification and accuracy.


\section{Outlook}
Future work should evaluate the systematic integration of additional spectral indices and channels to more accurately determine their influence on model performance.
Furthermore, the application and modification of additional \acrshort{YOLO} variants and comparison with Two-Stage detectors (e.g., \acrshort{R-CNN}-based methods) offer a promising field of investigation. 
Another approach is to analyse scaling effects, for example by reducing the spatial resolution to satellite level, in order to evaluate the transferability to raw data from real systems (e.g. Airbus, 30 cm/pixel \cite{airbus_neo}). 
The role of object-specific features, such as vehicle colour, should also be investigated in greater depth.
In the long term, the use of multitemporal data in particular opens up new perspectives, as it allows dynamic patterns and changes in the space-time context (e.g. commuter traffic versus long-term parking) to be captured.
